\chapter*{Preface: two changes of language}
\addcontentsline{toc}{chapter}{Preface: two changes of language}

These notes prove the Weil conjectures by making two changes of language
before beginning the arithmetic argument.  The first change is
coefficient-theoretic:
\[
  X_{\et}\text{ with towers }
  (\mathbf Z/\ell^n)_{n\geq1}
  \quad\rightsquigarrow\quad
  X_{\proet}\text{ with an honest }\Zl\text{-sheaf}.
\tag{P.1}\label{eq:preface-proetale-change}
\]
The second is cohomological:
\[
  \text{ordinary cohomological degree}
  \quad\rightsquigarrow\quad
  \text{middle perverse degree}.
\tag{P.2}\label{eq:preface-perverse-change}
\]
The first replaces the Artin--Rees quotient presentation in the proof spine
by a comparison-equivalent honest-sheaf working model.  The second makes
dimension, extension across a boundary, duality, and the Fourier transform
interact in the correct abelian category.
Together they give a cleaner proof and, more importantly for formalization, a
cleaner collection of reusable interfaces.

\begin{proetalebox}
On \(X_{\proet}\), the topological ring \(\Zl\) determines a genuine sheaf
\[
 \widehat{\mathbf Z}_{\ell,X}
   =\varprojlim_n\mathbf Z/\ell^n\mathbf Z.
\]
Thus coefficients and cohomology are operations on actual pro-etale
sheaves.  The naive category of all such sheaves is nevertheless too large:
proper base change fails there even for finite coefficients.  Route B first
passes to the reflective localization of solid pro-etale complexes,
\[
 \Dsolid(X,\Zl)\subset\Dnaive(X,\Zl),
\]
constructs the coherent six operations on that presentable ambient category,
and only then restricts them to \(\Dcons(X,\Zl)\).  Solidity is a continuity
condition, not a finite-stratification condition.  The historical Verdier
quotient of coefficient towers remains a comparison-equivalent model on
constructible objects.
\end{proetalebox}

\begin{powerbox}
Perverse sheaves are not merely used in the final three-line shift
calculation.  They perform three essential tasks in the proof:
\begin{enumerate}
\item affine Artin vanishing becomes one-sided perverse \(t\)-exactness;
\item Fourier--Deligne transform becomes an exact equivalence of perverse
      hearts;
\item the global weight theorem is expressed by the canonical objects
      \(\pH^a(Rf_!\cP)\), so boundary and support terms are controlled by
      subobjects, quotients, and intermediate extension.
\end{enumerate}
\end{powerbox}

\section*{What the pro-etale site simplifies}

Fix a prime \(\ell\) different from the characteristic of the ground field.
More generally, let \(E/\Ql\) be finite, let \(\mathcal O_E\) be its ring of
integers, and let \(\varpi\) be a uniformizer.  The coefficient ladder used in
the notes is
\[
\begin{tikzcd}[column sep=large]
 \widehat{\mathbf Z}_{\ell,X}
   \arrow[r] &
 \widehat{\mathcal O}_{E,X}
   \arrow[r,"\varpi^{-1}"] &
 E_X
   \arrow[r] &
 (\Qbarl)_X .
\end{tikzcd}
\tag{P.3}\label{eq:preface-coefficient-ladder}
\]
The last term is shorthand for descent to some finite \(E\), followed by
extension of scalars.  It is not treated as a complete noetherian coefficient
ring.

For integral coefficients, the Route B ambient category is
\(\Dsolid(X,\widehat{\mathcal O}_E)\).  Inside it, the category used by the
arithmetic proof consists of constructible, derived
\(\varpi\)-complete complexes.  For such a complex,
\[
  K\simeq
  \Rlim_n\bigl(K\otimes^L_{\widehat{\mathcal O}_E}
                    \mathcal O_E/\varpi^n\bigr).
\tag{P.4}\label{eq:preface-derived-completion}
\]
Weakly contractible pro-etale affines make this formula behave as the
geometric intuition suggests.  On a topologically noetherian scheme,
constructibility is equivalently local perfectness along a finite
stratification.

Formula~\eqref{eq:preface-derived-completion} is a compatibility and
recognition theorem; it is not the definition of an ambient operation.
Solid rigidity constructs the six operations before constructibility is
imposed.  Finiteness and constructible duality are then used to prove that
those operations preserve the constructible subcategory.

This is a simplification of the categorical architecture, not a magical
proof of the geometric theorems.  Proper and smooth base change, generic
local acyclicity, constructibility of direct images, absolute purity,
Verdier duality, and the trace formula remain deep.  The known proof of
ambient solid base change still uses torsion constructible generators and
classical torsion base change; Route B removes the levelwise definition, not
the finite theorem's mathematical content.

\begin{auditbox}
For a separated finite-presentation map, the exceptional direct image
\(Rf_!\) is first constructed in the ambient solid category by the coherent
compactification formalism and then restricted to constructibles.  It is
generally \emph{not} the derived functor of an abelian sheaf-level \(f_!\).
Pullback on the constructible complete subcategory agrees with completed
pullback.  Neither completion alone nor the unrestricted naive pro-etale
category satisfies the required ambient axioms; the counterexample and the
solid correction are proved in \cref{chap:proetale-six}.
\end{auditbox}

\section*{Why the coefficients change during the proof}

The integral category supplies lattices, reduction maps, continuous
cohomology, and perfect complexes.  The arithmetic weight argument is carried
out after passing to \(E\) or \(\Qbarl\).  There are two reasons.

First, an Artin--Schreier character
\(\psi:\mathbf F_p\to E^\times\) may require enlarging \(E\) so that it
contains the relevant \(p\)-th roots of unity.  Its rank-one kernel naturally
has an \(\mathcal O_E\)-lattice, but need not be defined over \(\Zl\).

Second, the familiar middle perverse \(t\)-structure is Verdier self-dual over
a field.  Over the discrete valuation ring \(\mathcal O_E\), duality exchanges
the two integral perversities \(p\) and \(p^+\).  Consequently integral
intermediate extensions come in two variants, and an integral decomposition
theorem is false in general.  These notes develop enough of the integral
pair to make reduction and lattices honest, then rationalize before invoking
self-duality, pure semisimplicity, or decomposition.

\section*{The perverse proof engine}

Let \(f_0:X_0\to Y_0\) be separated of finite type over \(\Fq\).  The central
theorem is first stated for an \(E\)-perverse Weil sheaf:
\[
 \cP_0\text{ has weights }\leq w
 \quad\Longrightarrow\quad
 \pH^a(Rf_{0!}\cP_0)
 \text{ has weights }\leq w+a.
\tag{P.5}\label{eq:preface-perverse-direct-image}
\]
The derived version follows from the perverse truncation tower.  This order
matters: perversity is the proof language of the theorem, not a restatement
attached after an ordinary-cohomology proof.

The difficult step is the relative curve estimate.  For a pure lisse sheaf
\(\cL\) on a smooth open curve \(j:U\hookrightarrow C\), one forms
\[
  \cP=j_{!*}(\cL[1]).
\tag{P.6}\label{eq:preface-curve-perverse-object}
\]
The Fourier--Deligne transform, normalized by a shift \([1]\), is exact on
the perverse heart.  Its stalk traces are exponential sums, and Plancherel
turns their second moment into a geometric norm.  Local Fourier transform
records the boundary contribution.  The resulting stalk estimate gives the
weight bound for \(H_c^1(U,\cL)\).  Relative Noether normalization,
recollement, and affine one-sided exactness then promote the curve result to
\eqref{eq:preface-perverse-direct-image}.

The complete noncircular chain is
\[
\begin{tikzcd}[column sep=small,row sep=large]
\substack{\text{pro-etale coefficients}\\
          \text{and six operations}}
 \arrow[r] &
\substack{\text{perverse recollement}\\
          \text{and Artin vanishing}}
 \arrow[r] &
\substack{\text{Fourier exactness}\\
          \text{and Plancherel}}
 \arrow[d] \\
\substack{\text{proper purity}\\
          \text{and IC purity}}
 \arrow[d] &
\substack{Rf_!\text{ upper}\\
          \text{weights}}
 \arrow[l] &
\substack{\text{perverse curve}\\
          \text{estimate}}
 \arrow[l] \\
\substack{\text{pure perverse}\\
          \text{semisimplicity}}
 \arrow[r] &
\substack{\text{decomposition}\\
          \text{theorem}}
 &
\\[-0.4em]
\substack{\text{smooth proper}\\
          H^i\text{ pure of weight }i}
 \arrow[r] &
\substack{\text{trace, duality,}\\
          \text{weight separation}}
 \arrow[r] &
\text{Weil conjectures.}
\end{tikzcd}
\tag{P.7}\label{eq:preface-dependency-chain}
\]
The smooth-proper purity node is obtained directly from proper purity; it
does not depend on the decomposition node.  Both are consequences of the
same stronger perverse theorem.

\section*{The stronger theorem proved first}

The natural endpoint of the perverse argument is stronger than the classical
smooth statement.  If \(X_0\) is proper and pure dimensional, and if
\(\cL_0\) is a pure lisse sheaf on a smooth dense stratum, then
\[
  \IC_{X_0}(\cL_0)=j_{!*}(\cL_0[d])
\]
is pure, and its intersection cohomology is pure with the predicted shift.
For a proper map \(f_0\), every
\(\pH^a(Rf_{0*}\IC_{X_0}(\cL_0))\) is pure and becomes semisimple after
geometric base change.  After base change to \(\Fbar\) and after forgetting
the Weil structure, the derived direct image splits noncanonically into these
canonical perverse cohomology objects.  No Frobenius-equivariant splitting
over \(\Fq\) is asserted; the perverse cohomology objects themselves, unlike
a chosen splitting, remain canonical and carry their Weil structures.

For a smooth \(d\)-fold,
\(\IC_{X_0}(E)=E[d]\).  Taking the structural morphism
\(a_0:X_0\to\Spec\Fq\) gives
\[
\begin{aligned}
 E[d]\text{ is pure of weight }d
 &\Longrightarrow Ra_{0*}E[d]\text{ is pure of weight }d,\\
 &\Longrightarrow
 H^{i-d}(Ra_{0*}E[d])
   =H^i(X,E)\text{ is pure of weight }i.
\end{aligned}
\tag{P.8}\label{eq:preface-terminal-calculation}
\]
This is the Riemann-hypothesis part.  The trace formula gives rationality,
Poincare duality gives reciprocal roots and the functional equation, and
universal-realization purity separates the rational zeta function into
integral, \(\ell\)-independent factors.

\section*{Completeness contract}

“Self-contained” in these notes has the following exact meaning.

\begin{enumerate}
\item Every coefficient category, functor, Weil structure, shift, twist,
      weight convention, perverse truncation, and comparison used in the
      proof is defined.
\item The implications from the foundational theorem dossier to
      \eqref{eq:preface-terminal-calculation} and then to every clause of the
      Weil conjectures are proved.
\item SGA-scale construction theorems are never hidden under the phrase
      “the usual formalism.”  Their exact hypotheses, compatibilities, and
      primary locations are recorded in the appendices.
\item The pro-etale/classical-adic equivalence is an explicit comparison
      theorem.  It is not treated as definitional equality.
\item No smooth-proper purity, decomposition theorem, integral Frobenius
      factor, or \(\ell\)-independence statement is used upstream of its
      proof.
\end{enumerate}

The foundational dossier distinguishes four statuses:

\begin{center}
\begin{tabular}{cl}
\toprule
\(\mathbf F\) & deep foundational theorem used with an exact source and
                interface,\\
\(\mathbf C\) & construction or comparison proved in the pro-etale chapters,\\
\(\mathbf P\) & theorem proved in full in the mathematical proof spine,\\
\(\mathbf A\) & final assembly from already proved results.\\
\bottomrule
\end{tabular}
\end{center}

These letters are expository metadata only.  They never occur as hypotheses
of the final theorem.

\section*{How to read the notes}

Readers interested in the new coefficient architecture should begin with
\cref{chap:proetale-site,chap:proetale-coefficients,chap:proetale-six,%
chap:proetale-weil}.  Readers who know the six operations may begin with the
middle perverse \(t\)-structure and the integral-lattice warning, then read
the perverse-engine and Fourier chapters.  The shortest logically complete
route to the final theorem is:
\[
 \text{perverse direct-image theorem}
 \longrightarrow
 \text{proper purity}
 \longrightarrow
 \text{terminal calculation}
 \longrightarrow
 \text{algebraic endgame}.
\]

The principal primary sources are Bhatt--Scholze and
Hemo--Richarz--Scholbach for pro-etale constructibles,
Ruimy--Tubach--Wolf \cite{RuimyTubachWolfSolid} for the solid Route B
ambient formalism, SGA and the finiteness/duality literature for the
geometric inputs, BBD \cite{BBD} for perverse sheaves,
Deligne \cite{DeligneWeilII} for weights, and Laumon
\cite{LaumonFourier} for the Fourier proof.  The exact theorem crosswalk is
given in the final appendix.
